\clearpage\begingroup\SourceLeaf{121}\HistoricalHead{116}{L}{PHYSIK. BRIEFWECHSEL.}
\begin{center}[30.]\end{center}
\begin{center}\textsc{Gauss} an \textsc{Bessel}.\quad Göttingen, 31. December 1831.\end{center}
\begin{center}\begin{minipage}{.88\linewidth}\centering\hrule\vspace{3pt}
\small [Briefwechsel zwischen \textsc{Gauss} und \textsc{Bessel}, Berlin 1880, S. 504.]
\par\vspace{3pt}\hrule\end{minipage}\end{center}
\noindent .\ .\ .\ .\ .\ .\ .\ .\ .\ .\ .\ .\ .\ .\ .\ .\ .\ .\ .\ .\ .\par
\indent Bei der Lectüre von \textsc{Poissons} älterer Abhandlung über die Electricität
(1812)\textsuperscript{[*)]} bin ich neulich auf einen meines Wissens neuen und wie mir scheint
artigen Satz gekommen. Er geht [a. a. O. S. 3] davon aus, dass die freie Electricität
sich unter der Oberfläche eines Isolirten Leiters sammelt, und meint, es
sei nicht zureichend, dass die Resultante aller Abstossungen in jedem Punkte
der innern Oberfläche jener Electricitätsschicht auf diese senkrecht sei, sondern
die physische Aufgabe fordere ausserdem noch die Bedingung, dass die Resultante
in jedem Punkte des innern eingeschlossenen Raums \(=0\) werde, weil sich
sonst neue Electricität zersetzen würde. Man mag diese Behauptung an sich
zugeben, allein ich \textls[120]{beweise}, dass, unter Voraussetzung der Repulsion im
verkehrten doppelten Verhältniss des Abstandes, die zweite Bedingung schon
von selbst in der ersten enthalten ist. Dann wird es offenbar ein Fehler, da
von neuem an die physische Natur zu appelliren, wo die Mathematik den
Gegenstand schon von selbst darreicht. Mein Satz gilt eben so gut von anziehenden
Kräften und heisst dann: Wenn ein wie immer gestalteter homogener
oder nicht homogener Körper aus Theilen besteht, die eine Anziehungskraft
im verkehrten doppelten Verhältniss des Abstandes ausüben, und dieser
Körper einen hohlen Raum umschliesst, an dessen Begrenzung in keinem
Punkte eine schiefe Resultante Statt hat, so ist die Resultante an jedem Punkte
dieser Begrenzung sowohl als in jedem Punkte des hohlen Raums \(=0\). —
Sobald man den Satz einmal aufgefasst hat, ist die Beweisführung eben nicht
schwer zu finden.\quad .\ .\ .\ .\ .\ .\ .\ .\ .\ .\ .\ .\ .\ .\ .\ .\ .\ .\ .\ .\ .
\begingroup\renewcommand{\thefootnote}{}\footnotetext{[*) S. D. \textsc{Poisson}, \textit{Mémoire sur la distribution de l'Electricité à la surface des corps conducteurs}, Mémoires de l'Institut, année 1811, Paris 1812, S. 1.]}\endgroup
\par\bigskip\begin{center}\rule{.24\linewidth}{.4pt}\end{center}
\clearpage\endgroup
